\documentclass[11pt,a4paper,sans]{moderncv}

\moderncvstyle{classic}
\moderncvcolor{blue}

\usepackage[scale=0.82]{geometry}

\newcommand{\doilink}[1]{\href{https://doi.org/#1}{doi:#1}}
\newcommand{\cvsubsection}[1]{\medskip\noindent\textbf{#1}\par\smallskip}

\firstname{Chris}
\familyname{Smith}
\address{International Institute for Applied Systems Analysis (IIASA)}{Schloßplatz 1, 2361 Laxenburg, Austria}
\email{chris@fairmodel.net}
\social[bluesky]{cjsmith.be}
\homepage{cjsmith.be}
\extrainfo{\href{https://orcid.org/0000-0003-0599-4633}{orcid} \textbar\ \href{https://scholar.google.com/citations?user=4P10oOEAAAAJ\&hl=en}{scholar}}

\begin{document}
\makecvtitle

I am a physical climate scientist with more than 10 years' experience, working at the interface of Earth System modelling and scenario projections. Presently, I am best known as the developer of the FaIR reduced-complexity climate model that was used extensively in the 2018 IPCC Special Report on Global Warming of 1.5°C, the IPCC Sixth Assessment (AR6) Working Group I (WG1) and Working Group III (WG3) reports, and several high-impact academic and policy outputs. I was a key contributor to the Reduced Model Intercomparison Project (RCMIP) that benchmarked the performance of emulators for use in the IPCC AR6. I am also a central part of the team that annually updates IPCC WG1 physical climate assessments based on new data (the Indicators of Global Climate Change project). I have been selected as a Coordinating Lead Author for the IPCC Seventh Assessment Report in WG1 Chapter 5 (Scenarios and Projected Future Global Temperatures).

The other major strand of my research is the Earth's energy budget and radiative forcing. I led and contributed to several important papers analysing the radiative forcing from CMIP6 era models, forming a critical component of the radiative forcing assessments in IPCC AR6 WG1. I am heavily involved in the design and delivery of the Coupled Model Intercomparison Project Phase 7 (CMIP7), having served as a member of the Strategic Ensemble Design Task Team, a current co-chair of the Radiative Forcing Model Intercomparison Project (RFMIP), and an advisory board member to the Scenario Model Intercomparison Project (ScenarioMIP).

\section{Career experience}
\cventry{2026--}{Senior Research Scholar}{Energy Climate and Environment Program, International Institute for Applied Systems Analysis}{Laxenburg, Austria}{}{}
\cventry{2024--2025}{Senior Scientist}{Department of Water \& Climate, Vrije Universiteit Brussel}{Brussels, Belgium}{}{}
\cventry{2023--2024}{Climate Mitigation Expert Scientist}{UK Met Office}{}{}{}
\cventry{2015--2024}{Senior Research Fellow}{School of Earth and Environment, University of Leeds}{}{}{}

\section{Prizes and recognition}
\cvline{2025}{Highlight talk, EGU 2025 \textit{Updated IPCC emissions scenarios no longer limit warming to 1.5°C}}
\cvline{2024}{Stanford/Elsevier Top 2\% Scientists 2024}
\cvline{2023}{Winner of the Societal Impact Award, UK Natural Environment Research Council}
\cvline{2018--2022}{Coordinating Lead Author, Lead Author, Chapter Scientist and Contributing Author in several roles across three reports of the Intergovernmental Panel on Climate Change (Sixth Assessment Report (AR6) Working Group 1, AR6 Working Group 3, and Special Report on Global Warming of 1.5°C)}
\cvline{2016}{Outstanding reviewer contribution award, Renewable Energy journal}

\section{Research funding}
\cvline{2024--2028}{Principal Investigator, HYWAY, Horizon Europe (consortium €4.6m)}
\cvline{2024--2027}{Work Package Leader, Climate Scenarios Compass Initiative (led by IIASA), Bezos Earth Fund (consortium €2.2m)}
\cvline{2024--2026}{Co-Investigator, MAGICA (led by U. Leeds), UK Natural Environment Research Council (consortium \pounds1m)}
\cvline{2024--2026}{Co-Investigator, Contrail assessment of future aircraft and propulsion architectures (led by U. Southampton), UK Natural Environment Research Council (consortium \pounds1m)}
\cvline{2023--2026}{Principal Investigator \& Work Package leader, WorldTrans, Horizon Europe (consortium €5m)}
\cvline{2023--2026}{Steering committee member, TRIFECTA, Norwegian Research Council}
\cvline{2020--2024}{Individual Fellowship, NERC-IIASA Collaborative Research Fellowship, Natural Environment Research Council (\pounds390k)}
\cvline{2017--2019}{Co-Investigator, CAMS74, Copernicus Atmosphere Monitoring Service radiative forcing products (\pounds101k)}

\section{Student and postdoctoral supervision}
\cvline{2023--}{Postdoctoral supervisor, Chris Wells, University of Leeds (WorldTrans)}
\cvline{2023--}{Lead PhD supervisor, Alejandro Romero Prieto, University of Leeds (PANORAMA Doctoral Training Programme, supported by Met Office CASE partner) \textit{Efficient models for climate mitigation and adaptation}}
\cvline{2023--}{PhD Co-supervisor, Magali Verkerk, University of Exeter \textit{Volcano-climate interactions}}
\cvline{2023--}{PhD Co-supervisor, Xinran Liu, University of Leeds \textit{Net-zero policy in China}}
\cvline{2023--2024}{MRes supervisor, Yongyao Liang, University of Leeds \textit{Attribution of 2023's record temperature}}
\cvline{2019--2024}{Seven MSc and undergraduate dissertation students at U. Leeds, two co-authored publications [32, 79]}

\section{Leadership and professional activities}
\cvline{2025}{Expert reviewer for Chancery Lane project for legal definitions of climate change}
\cvline{2024--2025}{Scenarios Forum 2025 (hosted by University of Leeds); part of winning bid team, and member of Scientific Steering Council and organizational committee}
\cvline{2023--2025}{External reviewer for three individual and team research proposals in Switzerland, Canada and the US}
\cvline{2023--}{Co-chair, Radiative Forcing Model Intercomparison Project (RFMIP) for CMIP7}
\cvline{2023--}{Committee member, CACTI (Composition Air-quality Climate inTeractions Initiative)}
\cvline{2023--2024}{External examiner for two PhD thesis defences in Canada and France}
\cvline{2022--2024}{Task Team Member, CMIP7 Strategic Ensemble Design Task Team}
\cvline{2021--}{External consultant on applications of reduced complexity climate models for high-profile clients including national governments}
\cvline{2021}{Panel member at COP26 IPCC side event, Glasgow, UK}
\cvline{2020--}{Convenor of European Geosciences Union sessions on climate modelling, climate emulation, economics, atmospheric chemistry, and air pollution}
\cvline{2019}{Scientific expert, British High Commission Malaysia Communicating Climate Science initiative, Malaysia}

\section{Media and outreach}
\cvline{2025}{Interview: Carbon Brief \textit{Implausibility of limiting warming to 1.5°C (forthcoming)}}
\cvline{2025}{Press conference: ``Hot takes \& policy quakes: When geoscience meets social science'', European Geosciences Union, Vienna, Austria \textit{On the feasibility of current scenarios to limit warming to 1.5°C}}
\cvline{2024}{Live national radio interview: BBC Radio Five Live \textit{On the climate ambitions of the two main political parties in the UK}}
\cvline{2024}{Press conference: United Nations Framework Convention on Climate Change (UNFCCC) SBSTA60, Bonn, Germany \textit{Indicators of Global Climate Change: Annual updates to IPCC climate system assessments}}
\cvline{2024}{TV interview: The Nation Kenya \textit{Solutions for climate change mitigation}}
\cvline{2024}{\href{https://www.youtube.com/watch?v=if9FR4VURqI}{YouTube interview}: British Embassy in Vienna \textit{Climate science with a message}}
\cvline{2023}{Live national radio interview: BBC Radio 4 Today Programme \textit{Northern Hemisphere 2023 heatwave}}
\cvline{2022}{Live TV interview: TVP Poland \textit{Climate impacts of NordStream gas pipeline leak}}
\cvline{2022}{TV interview: BBC Two Newsnight \textit{UK summer heatwave}}
\cvline{2018}{Live TV interview: Sky News \textit{IPCC Special Report on 1.5°C}}
\cvline{2018--}{Several guest posts on Carbon Brief and The Conversation climate-focused online news sources}
\cvline{2018--}{Several interviews on local radio and for newspapers in the UK, US, Germany, Austria, Australia, France, Norway, Sweden, Canada and Belgium. Examples include BBC, CNN, New York Times, Washington Post, The Guardian}

\section{Invited conference talks, workshops and panels}
\cvline{2025}{Belgian Ministry of Health: 10 years since the Paris Agreement}
\cvline{2025}{ESA Living Planet Symposium panel discussion, Vienna, Austria}
\cvline{2024}{Climate economics conference, Zürich, Switzerland}
\cvline{2024}{Econometric Models of Climate Change, Cambridge, UK}
\cvline{2024}{European Space Agency TRUTHS mission launch, Didcot, UK}
\cvline{2024--}{Transient Climate Response to Emissions workshops (Bristol, UK; Vienna, Austria)}
\cvline{2023}{Workshop on new modelling framework for climate mitigation, Paris, France}
\cvline{2023}{Gordon Research Conference on Radiation and Climate, Maine, USA}
\cvline{2023}{Climate emulator workshop, Vienna, Austria}
\cvline{2022}{Global Warming Levels workshop, Reading, UK}
\cvline{2022}{Institute for Mathematical and Statistical Innovation, Chicago, USA}
\cvline{2022}{Platform for Advanced Scientific Computing Conference, Basel, Switzerland}
\cvline{2020}{Institute of Physics, London, UK}
\cvline{2018}{American Geophysical Union Fall Meeting, Washington DC, USA}
\cvline{2017--}{PDRMIP, TriMIPathlon \& CACTI workshops (Imperial College, UK; IPSL, France; Princeton, US; Kiel, Germany; and online)}

\section{Invited institutional seminars}
\cvline{2025}{ETH Zürich, Switzerland}
\cvline{2025}{University of Oslo, Norway}
\cvline{2025}{Yonsei University, Seoul, South Korea}
\cvline{2024}{Vrije Universiteit Brussel, Belgium}
\cvline{2024}{Beijing Institute of Technology, China}
\cvline{2024}{University of Oslo, Norway}
\cvline{2022}{University of Zürich, Switzerland}
\cvline{2022}{BOKU, Vienna, Austria}
\cvline{2021}{University of Cambridge, UK}
\cvline{2020}{Energy Meteorology Group, University of Reading, UK}
\cvline{2019}{CICERO, Oslo, Norway}
\cvline{2019}{Imperial College, London, UK}
\cvline{2019}{University of East Anglia, Norwich, UK}
\cvline{2018}{Geophysical Fluid Dynamics Laboratory (GFDL), Princeton, New Jersey, USA}
\cvline{2018}{University of Manchester, UK}

\section{Visiting scientist positions}
\cvline{2025--}{Delft University of Technology, Delft, Netherlands}
\cvline{2025--}{Met Office Hadley Centre, Exeter, UK}
\cvline{2024--}{Priestley Centre for Climate Futures, University of Leeds, UK}
\cvline{2024}{CICERO, Oslo, Norway}
\cvline{2020--2025}{IIASA, Laxenburg, Austria}

\section{Teaching experience}
\cvline{2025}{Bergen Summer Research School, Bergen, Norway}
\cvline{2025}{Introduction to Physical Climate Change, UCLouvain, Belgium}
\cvline{2018--2024}{National Centre for Atmospheric Science PhD Introduction \textit{one-week intensive course for new PhD students from around the UK.} I delivered a lecture and practical session using a simple climate model demonstrating radiative forcing and the Earth's climate response.}
\cvline{2018--2021}{Climate Change \& Environmental Policy, University of Leeds, UK \textit{One-semester (22 teaching hours) course on the physical climate system aimed at social scientists.} I was module leader in 2020--21, and adapted the course substantially to maintain student interest and engagement due to shifting teaching activities online during the Covid-19 pandemic. Nominated by my students for an Inspiring Teaching Award in 2019.}
\cvline{2016--2021}{ad-hoc lecturing in the School of Earth \& Environment at University of Leeds, UK \textit{climate-related undergraduate courses.} Topics including climate mitigation scenarios, radiative forcing, the carbon cycle, geoengineering and sustainable food systems.}
\cvline{2013--2015}{Module demonstrator, School of Mechanical Engineering, University of Leeds, UK \textit{Thermofluids laboratory practicals}}

\section{Code and software development}
Proficient developer using collaborative software best practices, releasing version-controlled open-source software on the Python Package Index and Anaconda (fair and climateforcing packages). Experienced user of python, R, MATLAB, FORTRAN, bash, HTML, Git, GitHub, subversion, parallel computing, LaTeX and Unix systems. Experience in handling and processing large climate and observational datasets (netcdf4, hdf5). Experienced user of radiative transfer (SOCRATES, libRadtran, RFM) and complex climate (UKESM and HadGEM family) models, including code development.

\section{Education, qualifications, and previous career history}
\cventry{2011--2015}{Integrated MSc and PhD in Low Carbon Technologies}{School of Chemical and Process Engineering, University of Leeds}{UK}{}{}
\cventry{2008--2011}{International tax}{Deloitte}{Nottingham, UK}{}{}
\cventry{2004--2008}{Bachelor and Master of Mathematics (MMath)}{University of Nottingham}{UK}{}{}

\section{Publications}

\subsection{Reports}

\medskip\noindent\textbf{Intergovernmental Panel on Climate Change (IPCC)}\par\medskip

\textbf{\underline{Coordinating Lead Author}}
\begin{itemize}
\item \textbf{Seventh Assessment Report (AR7) Working Group 1 Chapter 5 (2025--2028) \textit{Scenarios and Projected Future Global Temperatures}}
\item Sixth Assessment Report (AR6) Working Group 1 Annex III, \textit{Tables of Historical and Projected Well-mixed Greenhouse Gas Mixing Ratios and Effective Radiative Forcing of All Climate Forcers} (Dentener, Hall, \textbf{Smith} et al., 2021)
\end{itemize}

\textbf{\underline{Lead Author}}
\begin{itemize}
\item AR6 Working Group 1 Chapter 7 Supplementary Material \textit{The Earth's Energy Budget, Climate Feedbacks and Climate Sensitivity} (\textbf{Smith} et al. 2021)
\end{itemize}

\textbf{\underline{Chapter Scientist}}
\begin{itemize}
\item AR6 Working Group 1 Chapter 7 (2021)
\item Special Report on 1.5°C Chapter 2 (2018)
\end{itemize}

\textbf{\underline{Contributing Author}}
\begin{itemize}
\item AR6 Working Group 3 Chapter 3 (2022)
\item AR6 Working Group 1 Summary for Policymakers (2021)
\item AR6 Working Group 1 Chapter 2 (2021)
\item AR6 Working Group 1 Chapter 4 (2021)
\item AR6 Working Group 1 Chapter 6 (2021)
\item AR6 Working Group 1 Chapter 7 (2021)
\item Special Report on 1.5°C Chapter 1 (2018)
\end{itemize}

\medskip\noindent\textbf{Other reports}\par\medskip
\begin{itemize}
\item \href{https://climate.leeds.ac.uk/wp-content/uploads/2024/06/Targets-for-effective-climate-mitigation-governance-in-the-UK-report-by-CEU.pdf}{Targets for effective climate mitigation governance in the UK}. Climate Evidence Unit, Leeds, 2024.
\item \href{https://www.lse.ac.uk/granthaminstitute/wp-content/uploads/2022/05/What-will-climate-change-cost-the-UK-risks-impacts-mitigation.pdf}{What will climate change cost the UK?} A study of climate risks, impacts and mitigation for the net-zero transition. London School of Economics, London, 2022.
\item ZERO IN ON series, produced under the CONSTRAIN Horizon Europe project 2019--2021 (\href{https://www.constrain-eu.org/wp-content/uploads/2021/11/Constrain-Zero-In-3-final.pdf}{vol. 3} \textbar\ \href{https://www.constrain-eu.org/wp-content/uploads/2020/02/CONSTRAIN-Zero-In-On-The-Remaining-Carbon-Budget-Decadal-Warming-Rates.pdf}{vol. 1})
\end{itemize}

\subsection{Peer-reviewed journal articles}

\input{pub_count.tex}
\pubcount\ total peer-reviewed, 5 in Nature/Science, h-index 50* (Google Scholar), > 12,000 citations*.

\url{https://orcid.org/0000-0003-0599-4633}

{\footnotesize *adjusted to account for peer-reviewed journal articles only}

\begingroup
\sloppy
\input{pubs_generated.tex}
\endgroup

\end{document}
